\documentclass[11pt,a4paper]{article}
\usepackage[a4paper,margin=2.5cm]{geometry}
\usepackage[spanish,es-nodecimaldot]{babel}
\usepackage{amsmath}
\usepackage{amssymb}
\usepackage{amsthm}
\usepackage{array}
\usepackage{booktabs}
\usepackage{longtable}
\usepackage{seqsplit}
\usepackage{hyperref}

\newtheorem{theorem}{Teorema}[section]
\newtheorem{lemma}[theorem]{Lema}
\newtheorem{proposition}[theorem]{Proposición}
\newtheorem{conjecture}[theorem]{Conjetura}
\theoremstyle{definition}
\newtheorem{definition}[theorem]{Definición}
\newtheorem{remark}[theorem]{Observación}

\title{Verificación computacional y análisis de redes de tensores 2D de los numeradores de politopos de orden de coronas multicapa}
\author{Wendel León Alvarado}
\date{\today}

\begin{document}
\maketitle

\begin{abstract}
Estudiamos los polinomios $h^*$ de Ehrhart de los politopos de orden asociados a los posets corona multicapa $C_{k,n}$. Una reducción propuesta a una traza de transferencia unidimensional es válida para altura dos, pero falla ya en $C_{3,3}$: para dilatación uno, el conteo exacto de puntos reticulares es $28$, mientras que la traza propuesta da $18$. Expresamos el conteo como una contracción de una red cilíndrica bidimensional de tensores con dimensión de enlace local $D=m+1$. La contracción sigue un orden jerárquico optimizado de operaciones tensoriales multilineales, sumando los índices compartidos según la convención de Einstein y sin truncar el rango de frontera; luego, el numerador de Ehrhart se recupera exactamente a partir de los conteos reticulares resolviendo un sistema triangular en la base binomial sobre $\mathbb{Z}$. Sucesiones racionales de Sturm certifican raíces reales, estrictamente negativas y simples en 14 casos, incluido $C_{7,4}$ de grado 21. La reconstrucción simbólica exacta sobre $\mathbb{Q}[t]$ produce recurrencias de orden $n-1$ para las circunferencias probadas $n=3,4$. La recurrencia universal de este orden se plantea como conjetura. Estos resultados constituyen evidencia finita, no una demostración para todo $k,n$.
\end{abstract}

\small\noindent\textbf{Palabras clave:} posets corona, polinomios de Ehrhart, redes de tensores, método de Sturm, raíces reales.
\normalsize

\section{Introducción y antecedentes}

Para un poset finito $P$, su politopo de orden se define por
\[
\mathcal{O}(P)=\{x\in[0,1]^P:x_u\le x_v\text{ siempre que }u\le_P v\}.
\]
Es un politopo reticular de dimensión $|P|$; sus vértices son los vectores característicos de los conjuntos superiores de $P$ \cite{Stanley1986}.

Fijemos $k\ge2$ y $n\ge3$. El poset $C_{k,n}$ tiene elementos $x_{j,i}$ para $1\le j\le k$ e $i\in\mathbb Z/n\mathbb Z$, con relaciones de cubierta
\[
x_{j,i}\prec x_{j+1,i},\qquad x_{j,i}\prec x_{j+1,i+1}\quad (1\le j<k),
\]
donde el segundo índice se interpreta cíclicamente. El poset tiene $kn$ elementos. Definimos
\[
L_{k,n}(m)=|m\mathcal{O}(C_{k,n})\cap\mathbb Z^{kn}|,\qquad
\sum_{m\ge0}L_{k,n}(m)q^m=\frac{H_{k,n}(q)}{(1-q)^{kn+1}}.
\]
Así, $H_{k,n}(q)=h^*(\mathcal{O}(C_{k,n}),q)$ y $H_{k,n}(0)=1$.

Lundström y Saud Maia Leite formularon la conjetura de raíces reales para la familia de coronas simples unidimensionales en trabajo iniciado en 2025 \cite{LundstromLeite2026}; Cao y Fu probaron posteriormente el caso de altura dos, $k=2$ \cite{CaoFu2026}. Este estudio computacional trata las familias multicapa $C_{k,n}$ con $k\ge3$. Sus aportes son el contraejemplo a la extensión propuesta de la traza unidimensional, una contracción de redes de tensores cilíndricas 2D, la reconstrucción exacta de numeradores de Ehrhart y el descubrimiento experimental de recurrencias. Las afirmaciones generales sobre recurrencias y ubicación de raíces siguen siendo conjeturas.

\subsection{Corrección metodológica: la traza unidimensional}

Para $N\ge1$, sea $K_N=(\min(a,b))_{1\le a,b\le N}$. En altura dos, con $N=m+1$, se tiene
\begin{equation}\label{eq:traza-altura-dos}
L_{2,n}(m)=\operatorname{tr}(K_N^n).
\end{equation}
En efecto, fijados los valores de la capa inferior, para valores adyacentes $a_{i-1},a_i$ el valor de su vecino común superior tiene $N-\max(a_{i-1},a_i)$ opciones. El cambio $a_i\mapsto N-a_i$ transforma ese factor en la entrada $\min(N-a_{i-1},N-a_i)$. Sumar alrededor del ciclo da \eqref{eq:traza-altura-dos}.

\subsection{Falla de la extensión superior propuesta}

Con tres capas o más, las variables intermedias permanecen acopladas por restricciones locales bidimensionales; sustituirlas por la matriz unidimensional propuesta no conserva el conteo. El siguiente ejemplo exacto refuta esa extensión.

\begin{proposition}\label{prop:contraejemplo}
Para $C_{3,3}$ y dilatación $m=1$, se tiene $L_{3,3}(1)=28$, mientras que la fórmula unidimensional propuesta da $\operatorname{tr}(K_2^3)=18$.
\end{proposition}

\begin{proof}
Para $m=1$, cada capa se representa mediante el subconjunto de sus tres posiciones cíclicas con valor $1$. Un subconjunto $A$ de una capa obliga a que la siguiente contenga $A\cup(A+1)$. El conteo de capas compatibles es $18+6+3+1=28$. En cambio,
\[
K_2=\begin{pmatrix}1&1\\1&2\end{pmatrix},
\qquad \operatorname{tr}(K_2^3)=18.
\]
Por tanto, la identidad de traza no es válida uniformemente para $k\ge3$.
\end{proof}

\section{Motor de redes de tensores 2D Boundary MPS}

Para una dilatación fija $m\ge0$, cada coordenada toma valores en $\{0,\ldots,m\}$. Definimos
\[
T_{a,b,c}=
\begin{cases}
1,&c\ge\max(a,b),\\
0,&\text{en otro caso},
\end{cases}
\qquad 0\le a,b,c\le m.
\]
Los tres índices tienen dimensión de enlace $D=m+1$. Los índices $a,b$ codifican valores adyacentes de una capa y $c$ el valor situado sobre ellos, imponiendo $x_{j+1,i}\ge\max(x_{j,i},x_{j,i+1})$. El conteo reticular es la contracción cilíndrica
\[
L_{k,n}(m)=
\sum_{\substack{x_{j,i}\in\{0,\ldots,m\}\\1\le j\le k,\ i\in\mathbb Z/n\mathbb Z}}
\prod_{j=1}^{k-1}\prod_{i\in\mathbb Z/n\mathbb Z}
T_{x_{j,i},x_{j,i+1},x_{j+1,i}}.
\]
La periodicidad en el índice circunferencial produce la geometría cilíndrica.

La red se contrae organizando sus factores multilineales a lo largo de la
frontera cilíndrica discreta y eligiendo una jerarquía eficiente de
contracciones de índices. La contracción abarca la red completa: no se
aplica compresión del rango de frontera ni truncamiento por valores
singulares. Así, «Boundary MPS» alude a la organización de la contracción
por fronteras, no a un método aproximado de estados matriciales.

La contracción exacta se efectúa sobre enteros de precisión arbitraria,
sin redondeo de punto flotante. Para la exploración espectral se considera
por separado una evaluación numérica con tensores reales o complejos de
doble precisión en el dispositivo de cómputo elegido. Estos valores
numéricos no se emplean como certificados enteros.

\section{Reconstrucción algebraica exacta y resultados de verificación}

\subsection{Reconstrucción exacta del numerador de Ehrhart}

Sea $d=kn$. La base binomial de Ehrhart satisface
\[
L_{k,n}(m)=\sum_{j=0}^{d}h_j\binom{m+d-j}{d},
\qquad H_{k,n}(q)=\sum_{j=0}^{d}h_jq^j.
\]
Evaluamos el conteo para $m=0,1,\ldots,d$ y resolvemos el sistema
triangular resultante mediante aritmética entera exacta. No se utiliza
interpolación de punto flotante.

\subsection{Resultados exactos}

La tabla presenta los 14 numeradores reconstruidos. El conteo de Sturm en $(-\infty,0)$ coincide con el grado y cada polinomio es libre de factores cuadrados; por ello todas sus raíces son reales, negativas y simples. Los coeficientes del polinomio se listan en orden descendente.

\small
\setlength{\tabcolsep}{3pt}
\begin{longtable}{@{}c c c >{\raggedright\arraybackslash}p{0.70\textwidth}@{}}
\caption{Dataset exacto y certificados de Sturm.}\label{tab:resultados}\\
\toprule
$k$ & $n$ & Grado; raíces negativas & Coeficientes de $H_{k,n}(q)$, en orden descendente\\
\midrule
\endfirsthead
\caption[]{Dataset exacto (continuación).}\\
\toprule
$k$ & $n$ & Grado; raíces negativas & Coeficientes en orden descendente\\
\midrule
\endhead
\bottomrule
\endlastfoot
2&3&4; 4/4&\texttt{\seqsplit{1,11,24,11,1}}\\
3&3&6; 6/6&\texttt{\seqsplit{1,18,93,160,93,18,1}}\\
4&3&8; 8/8&\texttt{\seqsplit{1,25,211,739,1120,739,211,25,1}}\\
5&3&10; 10/10&\texttt{\seqsplit{1,32,378,2088,5757,8064,5757,2088,378,32,1}}\\
6&3&12; 12/12&\texttt{\seqsplit{1,39,594,4550,19005,44547,59136,44547,19005,4550,594,39,1}}\\
7&3&14; 14/14&\texttt{\seqsplit{1,46,859,8468,48263,165086,344061,439296,344061,165086,48263,8468,859,46,1}}\\
8&3&16; 16/16&\texttt{\seqsplit{1,53,1173,14185,103331,473631,1393927,2657795,3294720,2657795,1393927,473631,103331,14185,1173,53,1}}\\
3&4&9; 9/9&\texttt{\seqsplit{1,69,1204,7164,16906,16906,7164,1204,69,1}}\\
4&4&12; 12/12&\texttt{\seqsplit{1,100,3058,37164,204639,553200,767068,553200,204639,37164,3058,100,1}}\\
5&4&15; 15/15&\texttt{\seqsplit{1,131,5869,116355,1144241,6007079,17690249,30135467,30135467,17690249,6007079,1144241,116355,5869,131,1}}\\
6&4&18; 18/18&\texttt{\seqsplit{1,162,9641,273808,4129580,35320040,179517436,560710448,1100896046,1376570220,1100896046,560710448,179517436,35320040,4129580,273808,9641,162,1}}\\
7&4&21; 21/21&\texttt{\seqsplit{1,193,14374,539238,11319081,141272409,1094268784,5432987952,17714213382,38591690550,56806162004,56806162004,38591690550,17714213382,5432987952,1094268784,141272409,11319081,539238,14374,193,1}}\\
3&5&12; 12/12&\texttt{\seqsplit{1,228,11096,175640,1128655,3314308,4717744,3314308,1128655,175640,11096,228,1}}\\
4&5&16; 16/16&\texttt{\seqsplit{1,349,32235,1090390,16336920,122193552,491902733,1115855805,1462319070,1115855805,491902733,122193552,16336920,1090390,32235,349,1}}\\
\end{longtable}
\normalsize

El caso límite $C_{7,4}$ tiene grado 21 y coeficiente central
$56\,806\,162\,004\approx5.68\times10^{10}$.

\subsection{Comprobaciones de sanidad y verificación de datos}

Los 14 polinomios certificados satisfacen
\[
H_{k,n}(0)=1,\qquad
q^{k(n-1)}H_{k,n}(1/q)=H_{k,n}(q),\qquad
\deg H_{k,n}=k(n-1).
\]
La segunda igualdad expresa la simetría palíndroma estricta (reciprocidad de Gorenstein); el grado coincide con el resultado estructural de la siguiente subsección. Además, cada certificado JSON registra un conteo de raíces negativas igual al grado y una comprobación de que el polinomio no tiene raíces múltiples. Estas verificaciones corresponden al dataset finito.

\subsection{Certificados racionales de Sturm}

Se construye la sucesión de Sturm de cada $H\in\mathbb Z[q]$ sobre $\mathbb Q$. Las variaciones de signo en $-\infty$ y en $0$ dan el número exacto de raíces distintas en $(-\infty,0)$. La simplicidad se comprueba por separado mediante $\gcd(H,H')=1$. Si el conteo coincide con $\deg H$, todas las raíces son reales, simples y estrictamente negativas.

La bisección racional permite aislar individualmente las raíces; cada
certificado JSON registra los coeficientes del polinomio y de la sucesión
de Sturm, las listas de signos, las variaciones y los resultados de
simplicidad y reciprocidad. Los 14 certificados exactos están
estructurados en formato JSON, nombrados
\texttt{certificado\_k\_n.json}, y disponibles en el repositorio
suplementario del proyecto.

\subsection{Grado y reciprocidad}\label{sec:estructura}

\begin{lemma}\label{lem:estructura}
Para $k\ge2$ y $n\ge3$, $H_{k,n}(q)$ es mónico, tiene grado $kn-k=k(n-1)$ y satisface
\[
H_{k,n}(q)=q^{kn-k}H_{k,n}(1/q).
\]
\end{lemma}

\begin{proof}
Los puntos reticulares interiores del múltiplo $m\mathcal O(C_{k,n})$ son aplicaciones estrictamente crecientes sobre las cubiertas con valores en $\{1,\ldots,m-1\}$. Cada cadena maximal tiene $k$ elementos, así que el menor múltiplo con puntos interiores es $m=k+1$. En ese múltiplo, las cadenas verticales fuerzan los valores $1,\ldots,k$, que satisfacen todas las cubiertas; hay exactamente un punto interior. Por la relación grado-código de Ehrhart, el grado es $kn+1-(k+1)=kn-k$, y el coeficiente principal es uno.

Para la reciprocidad, sea $\rho(x_{j,i})=j-1$. A cada aplicación interior $f$ del múltiplo $m$ se le asocia $g(x)=f(x)-\rho(x)-1$, una aplicación débilmente creciente con valores en $\{0,\ldots,m-k-1\}$. El cambio es biyectivo, de donde el conteo interior satisface $E^\circ(m)=E(m-k-1)$. Al combinar esta identidad con la reciprocidad de Ehrhart se obtiene la igualdad indicada.
\end{proof}

\section{Reconstrucción simbólica de recurrencias}

La resolución exacta de sistemas de coeficientes sobre $\mathbb Q[q]$ produce las siguientes proposiciones de alcance finito.

\begin{proposition}[Recurrencia finita exacta para $n=3$]
\label{prop:rec-n3}
Para las transiciones certificadas $k=4,\ldots,8$,
\[
H_{k,3}(q)=(q^2+6q+1)H_{k-1,3}(q)
+(q^3-2q^2+q)H_{k-2,3}(q).
\]
La solución es única y sobredeterminada dentro del espacio de búsqueda de grados acotados del analizador.
\end{proposition}

\begin{proposition}[Recurrencia finita exacta para $n=4$]
\label{prop:rec-n4}
Para las transiciones certificadas $k=6,7$,
\[
H_{k,4}(q)=P_4(q)H_{k-1,4}(q)+Q_4(q)H_{k-2,4}(q)+R_4(q)H_{k-3,4}(q),
\]
donde
\begin{align*}
P_4(q)&=q^3+23q^2+23q+1,\\
Q_4(q)&=7q^5-28q^4-22q^3-28q^2+7q,\\
R_4(q)&=q^8-5q^7+9q^6-5q^5\\
&\qquad{}-5q^4+9q^3-5q^2+q.
\end{align*}
La solución integral es única dentro del ansatz recíproco, con cotas de grados $(4,7,10)$ y spans recíprocos $(3,6,9)$.
\end{proposition}

Los operadores se reconstruyeron de manera única sobre $\mathbb Q[q]$ y se certificaron con aritmética entera exacta. La unicidad es relativa a los grados acotados para $n=3$ y, para $n=4$, a las cotas y spans recíprocos indicados; no se afirma unicidad sin esas restricciones. Los coeficientes son palíndromos respecto de sus spans estructurales. El orden observado es $n-1$ para $n=3,4$, pero las verificaciones abarcan solo las transiciones certificadas.

\section{Implicaciones para raíces reales y entrelazamiento}

\begin{conjecture}[Recurrencia cilíndrica universal de orden $n-1$]
\label{conj:recurrencia}
Para cada circunferencia fija $n\ge3$ existen polinomios
$A_{1,n}(q),\ldots,A_{n-1,n}(q)\in\mathbb Z[q]$, independientes de $k$,
tales que para todo $k\ge n$,
\[
H_{k,n}(q)=\sum_{j=1}^{n-1}A_{j,n}(q)H_{k-j,n}(q).
\]
\end{conjecture}

\begin{conjecture}[Raíces reales negativas en capas superiores]
\label{conj:raices}
Para todo $k\ge3$ y $n\ge3$, todas las raíces de $H_{k,n}(q)$ son reales, negativas y simples.
\end{conjecture}

\subsection{Discriminante característico para la recurrencia verificada de $n=3$}

Con la convención verificada $H_{k,3}=P_3H_{k-1,3}+Q_3H_{k-2,3}$,
$P_3(q)=q^2+6q+1$ y $Q_3(q)=q^3-2q^2+q$, la ecuación característica en $k$ es $\lambda^2-P_3(q)\lambda-Q_3(q)=0$. Por tanto,
\[
\Delta_3(q)=P_3(q)^2+4Q_3(q)
=q^4+16q^3+30q^2+16q+1.
\]
Para $q=-x<0$, sea $z=x+x^{-1}\ge2$. Entonces
\[
\frac{\Delta_3(-x)}{x^2}=z^2-16z+28=(z-2)(z-14).
\]
Así, el discriminante se anula en $z=2,14$, es positivo si $z>14$ y negativo para $2<z<14$. En consecuencia, no es positivo para todo $q<0$ y no implica por sí mismo estabilidad ni raíces reales para toda la familia. En particular, la expresión alternativa $q^4+16q^3+46q^2+8q+1$ vale $-63$ en $q=-4$. Los signos de los operadores aquí presentados son los que reproducen exactamente los polinomios certificados.

\subsection{Estado del entrelazamiento}

Los certificados de Sturm prueban que cada polinomio listado tiene raíces reales negativas simples; no prueban entrelazamiento estricto entre distintos valores de $k$. Los grados aumentan en $n-1$ al incrementar $k$ en uno, por lo que cualquier afirmación de entrelazamiento requiere un criterio generalizado explícito y una demostración independiente.

\section{Conclusiones y trabajo futuro}

La traza de transferencia de altura dos se obtiene al eliminar una sola capa, pero el conteo de $C_{3,3}$ muestra que la extensión unidimensional propuesta no conserva los conteos de puntos reticulares para capas superiores. La red de tensores cilíndrica preserva el acoplamiento local completo. La contracción exacta y la reconstrucción binomial producen polinomios reproducibles; los certificados de Sturm confirman raíces reales negativas simples para los 14 casos de la tabla. Las recurrencias solo se verificaron en los rangos finitos indicados; tanto su generalización universal de orden $n-1$ como la conjetura de raíces reales permanecen abiertas. Trabajo futuro incluye derivar las recurrencias de la estructura de transferencia y probar resultados uniformes de ubicación de raíces o entrelazamiento.

\subsection{Disponibilidad de datos y código}

El subdirectorio \texttt{ARQUITECTURA CON TENSORES/} contiene el motor
de redes de tensores y los módulos de reconstrucción exacta en \path{core/},
la interfaz de línea de comandos \path{main.py} y los 14 certificados JSON
en \path{certificados_json/}. La verificación de Sturm está implementada en
\path{core/root_verifier.py}; el análisis de recurrencias y entrelazamiento,
en \path{core/interlacing_analyzer.py}. Los certificados registran los
coeficientes y los datos exactos de Sturm necesarios para reproducir las
comprobaciones finitas.

\begin{thebibliography}{9}
\bibitem{LundstromLeite2026}
T.~Lundstr\"om y L.~Saud Maia Leite,
\emph{Order polytopes of crown posets},
European Journal of Combinatorics \textbf{133} (2026), 104304.
\href{https://doi.org/10.1016/j.ejc.2025.104304}{doi:10.1016/j.ejc.2025.104304}.

\bibitem{Stanley1986}
R.~P.~Stanley,
\emph{Two poset polytopes},
Discrete \& Computational Geometry \textbf{1} (1986), 9--23.

\bibitem{CaoFu2026}
Y.~Cao y H.~Fu,
\emph{Real-rooted Ehrhart $h^*$-polynomials of crown posets and wheel and whirl matroids},
preprint (2026), arXiv:2610.40496.
\href{https://arxiv.org/abs/2610.40496}{arXiv:2610.40496}.
\end{thebibliography}

\end{document}
